\section{意义是什么：从词到句子}

\subsection{意义并不是二元开关}

最后一讲最有哲学味道的部分，是对“meaning”的讨论。Manning 明确表达了一个立场：意义不是非黑即白，而是渐变的；它既来自真实世界的指称，也来自词语在各种使用场景中的关联方式。

\begin{figure}[H]
\centering
\includegraphics[width=0.95\textwidth]{slides-images/slide-040.jpg}
\caption{“红苹果在桌子上”的传统逻辑语义学例子：句法结构可以递归地组合出意义\protect\footnotemark}
\end{figure}
\footnotetext{Slides 展示了经典的组合语义树。}

\begin{knowledgebox}{“意义是梯度的”}
一个词或短语的意义不是“有”或者“没有”。更合理的理解是：我们对一个表达的理解程度是渐变的。你对一个词知道得越多、接触的语境越广，你对它的意义理解就越丰富。
\end{knowledgebox}

\subsection{use theory of meaning}

Manning 在这一段里提到了一个接近 Wittgenstein 的想法：词的意义来自它在语言和生活中的使用，而不仅仅是对应某个静态对象。

\begin{itemize}
    \item 有些词的意义主要通过指称对象来理解；
    \item 但很多词更多是通过语境、惯用法、社会实践和话语功能来理解；
    \item 因此，单纯“查字典式”的语义观并不够。
\end{itemize}

这也解释了为什么语言模型在某些意义上确实可以学到很多“用法知识”。它们看到的大量文本本身就是一个巨大的使用语料库。

\subsection{神经模型和语义的关系}

一个争论很集中的问题是：神经语言模型到底有没有语义？不同观点并不完全一致，但 Manning 的态度比较清楚：

\begin{itemize}
    \item 神经模型不只是“符号索引器”，它们确实学到了有关世界和语言使用的结构；
    \item 但这不等于它们已经具有和人类完全相同的理解方式；
    \item 意义很可能不只存在于词表定义里，也不只存在于纯粹的外部世界指称里。
\end{itemize}

\begin{figure}[H]
\centering
\includegraphics[width=0.95\textwidth]{slides-images/slide-056.jpg}
\caption{AI 的隐藏成本：信息、偏见、劳动、能耗、版权、隐私和环境代价都应该进入我们对“能力”的理解\protect\footnotemark}
\end{figure}
\footnotetext{Slides 在讨论 AI 的同时也提醒了现实世界的代价。}

\subsection{本章小结}

“意义”不是一个单点，而是一组相关但不相同的问题：词义、句义、使用、指称、语境、社会实践和可操作的任务行为。CS224N 的最后一讲并没有给出终极定义，但给出了一个更成熟的态度：不要把语义简化成任何一个单一视角。

%% ========== 风险与未来 ==========

